% Part: set-theory
% Chapter: spine
% Section: foundation

\documentclass[../../../include/open-logic-section]{subfiles}

\begin{document}

\olfileid{sth}{spine}{foundation}

\olsection{ভিত্তি}

আমাদের কাজ \emph{প্রায়} শেষ, কিন্তু \emph{পুরোপুরি}
নয়: $\ZFminus$-এর কোনো কিছু থেকেই নিশ্চিত হয়
না যে \emph{প্রত্যেক} সেট কোনো $V_\alpha$-র
সদস্য, অর্থাৎ প্রত্যেক সেট কোনো পর্যায়ে গঠিত।

স্তরক্রমের বাইরে সেট আছে কি না, সে বিষয়ে
একটি সরল (গাণিতিক) অর্থে আমাদের \emph{মাথা
ঘামাতে} হয় না: এমন সেট থাকলে সেগুলিকে উপেক্ষা
করতে পারি। কিন্তু সেটের \emph{ধারণাটি} আমরা
সমর্থন করেছি এই ভেবে যে প্রত্যেক সেট কোনো
পর্যায়ে গঠিত হয় (\olref[z][story]{sec}-এর
\stageshier{} দেখুন)। তাই স্তরক্রমের বাইরে সেট
থাকার সম্ভাবনা বাদ দিতে চাই। সে জন্য এমন একটি
নতুন স্বতঃসিদ্ধ যোগ করতে হবে, যা প্রত্যেক সেটের
স্তরক্রমে কোথাও অবস্থান নিশ্চিত করে।

$V_\alpha$-গুলিই আমাদের পর্যায়; তাই নিচের
বক্তব্যটিকে সরাসরি স্বতঃসিদ্ধ করার কথা ভাবতে
পারি:

\begin{defish}
\emph{নিয়মিততা।} $\forall A \exists \alpha\, A \subseteq V_\alpha$
\end{defish}

এটি যথার্থ একটি পথ। তবে পরের বিভাগে আলোচিত
কারণে আমরা এর বদলে অন্য একটি স্বতঃসিদ্ধ গ্রহণ
করব:

\begin{axiom}[ভিত্তি]
$(\forall A \neq \emptyset)(\exists B \in A)A \cap B = \emptyset$।
\end{axiom}

কিছু পরিশ্রমে দেখানো যায় যে $\ZFminus$-এ
ভিত্তি থেকে নিয়মিততা অনুসৃত হয়:
\begin{defn}
প্রত্যেক সেট $A$-এর জন্য ধরি:
\begin{align*}
	\text{cl}_0(A) &= A,\\
	\text{cl}_{n+1}(A) &= \bigcup \text{cl}_n(A),\\
	\text{trcl}(A) &= \bigcup_{n < \omega} \text{cl}_{n}(A).
\end{align*}
$\text{trcl}(A)$-কে $A$-এর \emph{পরিযায়ী আবরণ}
বলি।
\end{defn}
\noindent
``পরিযায়ী আবরণ'' নামটি যথার্থ:
\begin{prop}\ollabel{subsetoftrcl}
$A \subseteq \trcl{A}$ এবং $\trcl{A}$ একটি
পরিযায়ী সেট।
\end{prop}

\begin{proof}
স্পষ্টতই $A = \text{cl}_0(A) \subseteq \text{trcl}(A)$।
আর $x
\in b \in \trcl{A}$ হলে কোনো $n$-এর জন্য
$b \in \text{cl}_n(A)$; তাই $x
\in \text{cl}_{n+1}(A) \subseteq \text{trcl}(A)$।
\end{proof}

\begin{lem}\ollabel{lem:TransitiveWellFounded}
$A$ পরিযায়ী সেট হলে এমন একটি $\alpha$ আছে
যে $A
\subseteq V_\alpha$।
\end{lem}

\begin{proof}
\olref[ordinals][opps]{defsupstrict}-এ
``$\supstrict(X)$''-এর সংজ্ঞা মনে রেখে দুটি
সেট সংজ্ঞায়িত করি:
\begin{align*}
	D  &= \Setabs{x \in A}{\forall \delta\ x \nsubseteq V_\delta}\\
	\alpha &= \supstrict\Setabs{\delta}{(\exists x \in A)
	(x \subseteq V_\delta \land (\forall \gamma \in \delta)x \nsubseteq V_\gamma)}
\end{align*}
% BN-SRC-789 TARGET-CORRECTION-BEGIN
\emph{উৎস-সংশোধন-নোট।} এখানে পর্যায়ের সূচক হিসেবে ব্যবহৃত
গ্রিক চলগুলি ক্রমসংখ্যায় বিস্তৃত; নিয়মিততার $\alpha$-ও তাই।
$D$-র জন্য পৃথকীকরণ যথেষ্ট। কিন্তু $\alpha$-র সংজ্ঞায়
প্রতিস্থাপন আগে $A\setminus D$-তে প্রয়োগ করি: ঠিক সেই সদস্যগুলির
কোনো পর্যায়ে উপসেট হিসেবে অবস্থান আছে। সূত্রভিত্তিক সর্বনিম্নতার
নীতিতে প্রত্যেকটির একমাত্র ক্ষুদ্রতম $\delta$ আছে, তাই তাদের
image একটি সেট। $D$ শূন্য হওয়ার আগেই এই bounded domain
সংজ্ঞাটি বৈধ; ভালো সদস্য না-থাকলে image শূন্য।
আর $x\subseteq V_\delta$ ও $\delta\in\alpha$ থেকে
$x\in V_{\delta+1}\subseteq V_\alpha$: ঘাত সেটের সংজ্ঞা
এবং $\delta+1\leq\alpha$ ব্যবহার করা হয়েছে। পরের $\beta$-র
image-এ একই যুক্তি গোটা $B$-তে চলে, কারণ সেখানে প্রত্যেক
সদস্যই $D$-র বাইরে।
% BN-SRC-789 TARGET-CORRECTION-END
ধরি, $D = \emptyset$। তাহলে $x \in A$ হলে
এমন একটি $\delta$ আছে যে $x \subseteq V_\delta$;
ক্রমসংখ্যাগুলির সুক্রমের কারণে সেই ক্রমসংখ্যাটি
ক্ষুদ্রতম নেওয়া যায়, অর্থাৎ
$(\forall \gamma \in \delta)x \nsubseteq V_\gamma$।
তাই $\delta \in \alpha$ এবং
\olref[spine][Valphabasic]{Valphabasicprops}
অনুযায়ী $x \in V_\alpha$। সুতরাং চাওয়া মতো
$A \subseteq
V_{\alpha}$।

অতএব $D = \emptyset$ দেখালেই যথেষ্ট। বিরোধের
উদ্দেশ্যে অন্যথা ধরি। ভিত্তি-স্বতঃসিদ্ধে
এমন একটি $B \in D \subseteq A$ আছে যে
$D \cap B
= \emptyset$। $x \in B$ হলে $A$-র পরিযায়িতায়
$x \in A$; আর $x \notin D$ বলে
$\exists \delta\ x \subseteq
V_\delta$। তাই এবার ধরি
% BN-SRC-440: উৎসে মুক্ত ছোটহাতের b ছিল; ভিত্তি-স্বতঃসিদ্ধে বেছে নেওয়া সেটটি B।
\[
	\beta = \supstrict\Setabs{\delta}{(\exists x \in B)(x \subseteq V_\delta \land (\forall \gamma < \delta)x \nsubseteq V_\gamma)}.
\]
আগের মতো $B \subseteq V_\beta$ পাই, যা
$B \in
D$-র বিরোধী।
\end{proof}

\begin{thm}\ollabel{zfentailsregularity}
নিয়মিততা সত্য।
\end{thm}

\begin{proof}
যেকোনো $A$ নিই। \olref{subsetoftrcl}
অনুযায়ী $A \subseteq \trcl{A}$ এবং
এই আবরণ পরিযায়ী। তাই
\olref{lem:TransitiveWellFounded} থেকে
এমন একটি $\alpha$ আছে যে
$A \subseteq \trcl{A} \subseteq V_\alpha$।
\end{proof}

এই ফলগুলি দেখায়, $\ZFminus$-এ
$\emph{ভিত্তি}\Rightarrow\emph{নিয়মিততা}$
প্রমাণ করা যায়।
\olref[spine][rank]{zfminusregularityfoundation}-এ
দেখাব, $\ZFminus$-এ
$\emph{নিয়মিততা}\Rightarrow\emph{ভিত্তি}$-ও
প্রমাণ করা যায়। তাই $\ZFminus$-এর সাপেক্ষে
ভিত্তি ও নিয়মিততা \emph{তুল্য}। অর্থাৎ
$\ZFminus$ ধরে নিয়মিততার সঙ্গে সমতুল্য বলেই
ভিত্তির ন্যায্যতা দিতে পারি; আর
\stageshier{} থেকে নিয়মিততাকে সরাসরি
ন্যায্যতা দিতে পারি।

\end{document}
